% =====================================================================
%  GMN3 §2  BPS 线算子（BPS line operators）
%  原文：Lamination-Paper.tex 561--1028 行；PDF p.7--13（印刷 p.6--12）
%  编号公式 (2.1)--(2.15) 共 15 条，全部在本块内，故 \eqref 可用；
%  跨块引用（节号、式号、附录）暂写原文编号并加 % TODO，合稿时统一改 \ref/\eqref
% =====================================================================

\section{BPS 线算子}
\label{sec:Susy-BPS-line}

\subsection{$\zeta$ 型线算子的定义}

本文沿用 \cite{Kapustin:2005py,Kapustin:2006pk} 所主张的线算子定义。我们考虑的是由紫外不动点理论经相关微扰而得到的量子场论；所谓线算子，就是该紫外理论在 $AdS_2\times S^2$ 上的一个共形不变边界条件。

本文聚焦于具有 $d=4$、$\CN=2$ 超对称的理论。因此其紫外不动点具有 $su(2,2\vert 2)$ 超共形对称性。（我们的约定见附录 A% TODO 合稿改 \ref{app:Susy-Conventions}
。）我们通常考虑 $\IR^{1,3}$ 中位于空间原点 $x^1=x^2=x^3=0$、沿时间方向伸展的线算子。此时，共形代数 $so(2,4)$ 的未破缺子代数由 $D,P_0,K_0$（它们生成 $sl(2,\IR)$ 的一个拷贝）与空间转动 $M_{ij}$（它们生成 $so(3)$ 的一个拷贝）生成。于是，定义我们线算子的那个紫外边界条件保持 $sl(2,\IR)\oplus so(3)\cong so(4^*)$，即 $AdS_2\times S^2$ 的等距变换李代数。此外，我们只研究保持超共形代数之 $su(2)_R$ R 对称性的算子。我们还要求保持一半超对称，因此要研究保持超代数 $osp(4^*\vert 2)$ 的线算子。事实上，选取该子代数的奇生成元有许多种方式，而这个族的存在在后续讨论中颇为重要。超代数上有一个对合 $I_\zeta$，它由超共形群的某个 $U(1)_R$ 变换共轭空间反射 $x^i\to -x^i$ 而诱导。在庞加莱超对称上，该对合的作用为
\begin{equation}
\begin{split}
Q_\alpha^A & \rightarrow \zeta^{-1}  \epsilon^{AB}
\sigma^0_{\alpha\dot \beta}\bar
Q^{\dot \beta}_{~ B} \\
\bar Q^{\dot \alpha}_{~ A} & \rightarrow  \zeta  \epsilon_{AB}
\bar\sigma^{0 \dot \alpha\beta} Q_\beta^{~B}
\end{split}
\end{equation}
而对特殊共形超对称 $S_\alpha^{~A},\bar S^{\dot \alpha A}$ 有类似的作用。这里 $\zeta$ 是由 $U(1)_R$ 转动得到的相位。该自同构的不动子代数记作 $osp(4^*\vert 2)_\zeta$；对一切 $|\zeta|=1$，它同构于 $osp(4^*\vert 2)$。事实上，所有这样的子代数都由 $U(1)_R$ 变换互相转过去。我们把保持上述超代数的线算子记作 $L(\zeta;\cdots)$ 或 $L_\zeta(\cdots)$，其中 $\cdots$ 指刻画某个具体线算子的其他参数，如下文第 \ref{subsec:LabelsForLines} 节所讨论的那些。正如第 4 节% TODO 合稿改 \ref{sec:Formal-line-operators}
与第 10 节% TODO 合稿改 \ref{sec:Examples}
的例子将显示的，线算子一般只在 $\zeta$ 圆周的万有覆盖上才单值。

我们要研究的场论，由各自的超共形不动点出发、通过沿库仑分支移动或通过开启质量而微扰得到。在理论的红外描述中有 $d=4$、$\CN=2$ 的庞加莱超对称。对合 $I_\zeta$ 也作用在这个超代数上，并附加一条规则 $I_\zeta: Z\to\zeta^2\bar Z$，其中 $Z$ 是中心荷算子。为方便起见，我们取 $\zeta^{-1}$ 的一个平方根，记作 $\xi$，并用下式定义在该对合下不变的超荷：
\begin{equation}\label{eq:Q-zeta-susy}
\CR_{\alpha}^{~A} = \xi^{-1} Q_{\alpha}^{~A} + \xi \sigma^0_{\alpha
\dot \beta } \bar Q^{\dot \beta A}
\end{equation}

沿时间方向伸展、且带线算子 $L$ 的理论，其希尔伯特空间记作 $\CH_L$。$\CR_\alpha^{~A}$ 就是作用在这个希尔伯特空间上的算子。

\subsection{线算子的半环}\label{subsubsec:RingStructure}

一个有趣的观察 \cite{Kapustin:2006hi}（在 \cite{Kapustin:2007wm} 中有更详细的展开）是：$\CN=2$ 规范理论中的超对称线算子构成一个半环（semiring），其结构相当精致，并细致地依赖于理论的物质内容。

为定义加法，我们说：两个线算子之和 $L+L'$ 的关联函数，就是 $L$ 与 $L'$ 各自关联函数之和。存在两个算子之和时，理论的希尔伯特空间就是两个超选择扇区之和，即分别与这两个算子相伴的两个希尔伯特空间：
\begin{equation}\label{eq:sum-line}
\CH_{L+L'}:= \CH_L \oplus \CH_{L'}.
\end{equation}
我们把不是另外两个线算子之和的线算子称为\emph{单线算子}。

此外，还可以通过在路径积分中插入两个线算子来定义算子的乘积 $LL'$。我们所考虑的线算子有一条特殊性质：它们的关联函数与两个算子的间距无关。为看清这一点，注意到那些在对合 $I_\zeta$ 下为奇的超对称 $T_{\alpha}^A = \xi^{-1} Q_\alpha^A -  \xi \sigma^0_{\alpha \dot \beta } \bar Q^{\dot \beta A}$ 满足对易关系
\begin{equation}
[\CR_\alpha^A, T_\beta^B ] = 2 \epsilon_{\alpha\beta} \epsilon^{AB}
(\bar   Z - \zeta^{-2} Z) + 8 \zeta^{-1} \epsilon^{AB}
\sigma^{0m}_{\alpha\beta} P_m.
\end{equation}
取 $\alpha\beta$ 的对称部分，我们就能把空间平移算子 $P_i$ 写成某个 $\CR_\alpha^A$ 的对易子；因此，所有沿 $x^0$ 轴对齐的型 $I_\zeta$ 线算子，其关联函数都与空间坐标无关。我们可以让这些算子互相靠近，而由局域性，其乘积应当等价于某个别的线缺陷。于是，算子乘积展开应当定义一种环乘法。下文我们将对若干有趣的例子把它确定下来。把两个线算子之乘积对应的希尔伯特空间与这两个线算子各自的希尔伯特空间联系起来，是一件有趣的事。这可以做，但出人意料地微妙。它将引导我们在第 3.6 节% TODO 合稿改 \ref{subsec:Deformed-Ring}
定义线算子乘积的一种有趣的非交换形变。

\subsection{拉氏规范理论中线算子的标签}\label{subsec:LabelsForLines}

许多 $d=4$、$\CN=2$ 理论都容许用矢量多重态与超多重态写出的拉氏量描述。对这类理论，按照 \cite{Kapustin:2005py} 与 \cite{Kapustin:2006hi}，我们可以对我们所关心的线算子的定义说得更具体（注意我们的 $\zeta$ 在后一文献中写作 $t^{-1}$）。

我们先回顾一些李代数的理论。从单李代数 $\lieg$ 出发，定义若干标准格。取嘉当子代数 $\liet\subset\lieg$。于是有一组典范的根 $\Phi(\lieg)\subset\liet^*$，它们来自把 $\liet$ 在 $\lieg$ 上的伴随作用对角化。对每个根 $\alpha\in\Phi(\lieg)$，都有 $\lieg$ 中一个 $sl(2)_\alpha$ 的拷贝；我们可以典范地定义相应的\emph{余根} $H_{\alpha}\subset\liet$，它生成 $sl(2)_\alpha$ 的嘉当子代数，并归一化为 $\langle \beta , H_{\alpha} \rangle = 2 (\beta,\alpha)/(\alpha,\alpha)$。（左端是 $\liet^*$ 与 $\liet$ 的典范配对，右端用了嘉当–基灵型；但因假定 $\lieg$ 单，归一化会消掉。）根与余根分别生成格 $\Lambda_r\subset \liet^*$ 与 $\Lambda_{cr}\subset \liet$。关于余根格，有两种有用的看法。由李定理，存在唯一的单连通连通紧李群 $\tilde G$，其李代数为 $\lieg$。余根格的元素可以等同于 $U(1)$ 到 $\tilde T$ 的同态，其中 $\tilde T$ 是 $\tilde G$ 的极大环面，李代数为 $\liet$。另一种说法是：$\Lambda_{cr}$ 是 $\exp[ 2\pi \cdot]: \liet \to \tilde T$ 的核。

利用 $\liet$ 与 $\liet^*$ 的完美配对（而\emph{不是}基灵型），我们便典范地定义了权格与磁权格：
\begin{equation}
\begin{split}
\Lambda_{wt} & := \Lambda_{cr}^* \subset \liet^* \\
\Lambda_{mw} & := \Lambda_{r}^* \subset \liet.
\end{split}
\end{equation}
一个标准结论说，中心 $Z(\tilde G)$ 同构于商
\begin{equation}
Z(\tilde G) \cong \Lambda_{mw}/\Lambda_{cr} \cong
\Lambda_{wt}/\Lambda_r.
\end{equation}
最后，设 $\CW$ 是李代数 $\lieg$ 的外尔群。它同构于 $N(\tilde T)/\tilde T$，其中 $N(\tilde T)$ 是 $\tilde T$ 在 $\tilde G$ 中的正规化子。

现在，我们的理论一般带物质场，因此应当考虑一般的紧连通单李群，其李代数为 $\lieg$。仍由李定理，这样的群必形如 $G\cong \tilde G/\CZ$，其中 $\CZ\subset Z(\tilde G)$ 是中心的一个子群。于是 $Z(G) \cong Z(\tilde G)/\CZ$。由于中心在外尔群作用下不动，我们可以找到一个 $\CW$ 不变的格 $\Lambda_G$，使得
\begin{equation}
 \Lambda_{cr}\subset \Lambda_G \subset \Lambda_{mw}
 \end{equation}
从而 $\CZ \cong  \Lambda_G/\Lambda_{cr} $。事实上，可以把 $\Lambda_G$ 等同于群 $G$ 的指数映射的核：
\begin{equation}
\Lambda_G = \{ P \in \liet \vert \exp(2\pi P) = 1_G \}
\end{equation}
也就是说 $\Lambda_G \cong {\rm Hom}(U(1), T)$，其中 $T= \tilde T/\CZ$ 是 $G$ 的极大环面。对偶地，$G$ 的特征标群是 ${\rm Hom}(T, U(1))$，它典范地对偶于 ${\rm Hom}(U(1),T)$，因为存在完美配对
\begin{equation}
{\rm Hom}(U(1),T) \times {\rm Hom}(T,U(1)) \rightarrow {\rm
Hom}(U(1),U(1))\cong \IZ
\end{equation}
于是特征标群同构于 $\Lambda_G^*$；正是为此，$\Lambda_G$ 有时被称为 $G$ 的\emph{余特征标格}。$\Lambda_G^*$ 可以等同于权格 $\Lambda_{wt}(G)$。这个格也可以看成 $\tilde G$ 的那些在 $\CZ$ 下平凡变换的表示中的权。以上便完成了我们对李代数理论的回顾。

如我们所说，有了显式的场多重态，就能构造一些显式的线算子。最显然的线算子，是沿时间方向一条路径 $p$、在某个固定 $\vec x$ 处的超对称 Wilson 线 \cite{Baulieu:1997nj,Maldacena:1998im,Rey:1998ik}，取规范群的一个表示 $\rho_\CR: G\to {\rm End}(V)$，则
\begin{equation}\label{eq:Susy-Wilson-Zeta}
L_\zeta(\vec x ;\CR) =
\rho_{\CR} P\exp \oint_{p} \left( \frac{\varphi}{2\zeta } - i A
- \frac{\zeta \bar \varphi}{2 }  \right).
\end{equation}

此外还可以构造 't Hooft 算子。回忆一下，它们由对线算子无穷小邻域中的场施加边界条件来定义。在我们的情形，我们要求场在环绕的小球面 $S^2$ 上的极限值对应于一个保持 $osp(4^*\vert 2)_\zeta$ 之庞加莱子代数的超对称磁单极子。回忆：规范群为 $G$ 的理论中的球对称磁单极子，由 $S^2$ 上一个 $G$-丛的转移函数指定（它决定了物理空间 $\IR^4-\{p\}$ 上的主丛）。这样的 $G$-丛等同于 $\Lambda_G \cong {\rm Hom}(U(1),T)$。事实上，对任何这样的同态，都可以通过略微修改 Kapustin 的构造，造出一个保持庞加莱子代数的解（见附录 C% TODO 合稿改 \ref{app:Fixed-Point-Equations}
的式 (C.3)% TODO 合稿改 \eqref{eq:Embed-fix-sol}
）。于是每个同态都定义一个 't Hooft 线算子 $L_\zeta (p;v)$，其中 $v \in \Lambda_G$。\footnote{这里重要的是把 $\zeta$ 理解为在 $U(1)_R$ 的覆盖空间中取值。这一点从我们回顾过的定义看不出来。然而，量子反常把相位 $\zeta$ 的改变与 $\theta$ 角的平移联系起来；线算子的威滕效应表明，'t Hooft 算子不可能是 $\zeta$ 的单值函数。}

众所周知，可以把上述两种构造结合起来，定义 Wilson–'t Hooft 算子 $L_\zeta(p;v,\CR)$，其中 $\CR$ 是 $v$ 的中心化子的一个表示。如 \cite{Kapustin:2005py} 所示，这样的 Wilson–'t Hooft 算子的集合\emph{先验地}以 $\Lambda_{mw} \times \Lambda_{wt}$ 的元素为标签。然而，相容的线算子集合上还有局域性条件。例如，在带有标签为 $(P,0)$ 的 't Hooft 算子时，标签为 $(0,Q)$ 的 Wilson 线算子只对规范丛在 $\IR^3- \{0 \}$ 上的结构群的真表示（而非射影表示）才有意义。更一般地，我们预期：相容的单线算子集合可以由极大子格 $\CL\subset \Lambda_{mw} \times  \Lambda_{wt}$ 构造，这些子格满足一个类 Dirac 的量子化条件：
\begin{equation}\label{eq:Line-Op-Locality}
\forall (P,Q), (P',Q')\in \CL : \qquad \qquad \langle P, Q' \rangle - \langle P' , Q \rangle \in \IZ
\end{equation}
由于相差外尔变换的算子彼此规范等价，我们应当要求 $\CL$ 是 $\CW$ 不变的；此时线算子的规范同构类与 $\bar\CL=\CL/\CW$ 一一对应。这样一族相容集合的自然例子是
\begin{equation}\label{eq:Wilson-tHooft-B}
\bar\CL_G := (\Lambda_G \times \Lambda_G^*)/\CW
\end{equation}
在两个极端，对 $G=\tilde G$ 我们有 $(\Lambda_{cr}\times \Lambda_{wt})/\CW$，而对伴随群 $G = G_{\rm ad}:= \tilde G/Z(\tilde G)$ 我们有 $(\Lambda_{mw} \times \Lambda_{r})/\CW$。

最一般的极大互局部子格 $\CL$ 处于正合列
\begin{equation}\label{eq:label-seq}
0 \rightarrow \Lambda_G^* \rightarrow \CL \rightarrow \Lambda_G  \rightarrow 0
\end{equation}
之中。为看出这一点，先考虑投影 $(P,Q) \to P$。它定义了 $\Lambda_{mw}$ 的一个子格，而后者对某个 $G$ 必为 $\Lambda_G$。又因 $\CL$ 极大且互局部，该映射的核必形如 $(0,Q)$，其中 $Q \in \Lambda_G^*$。序列 \eqref{eq:label-seq} 虽然分裂，但并非\emph{自然}分裂——这反映了威滕效应的可能性——所以若考虑局部系统（例如在耦合常数空间上、或在 $\zeta$ 平面 $\IC^*$ 上），\eqref{eq:label-seq} 可以是非平凡的。

对线算子还可以陈述一个更一般的互局部条件，它应当适用于非拉氏理论。存在线算子 $L$ 时的希尔伯特空间 $\CH_L$ 应当是空间转动群 $SU(2)$ 的一个表示。若考虑位于 $\vec x_1, \vec x_2\in \IR^3$ 的两个线算子 $L_1,L_2$，则转动群破缺为群 $U(1)_{\vec x_1 \vec x_2}$，后者双重覆盖绕 $\vec x_1 \vec x_2$ 轴的转动群。与两线算子之乘积相伴的希尔伯特空间 $\CH_{L_1 L_2}$ 当然是 $U(1)_{\vec x_1 \vec x_2}$ 之万有覆盖的表示。我们的互局部条件就是如下陈述：$\CH_{L_1 L_2}$ 是 $U(1)_{\vec x_1 \vec x_2}$ 的真表示。在具有拉氏量描述的 $\CN=2$ 规范理论中，只要沿库仑分支走得足够远、使由线算子源出的场可以放进嘉当子代数，就能回到前面关于互局部条件的表述。这时可以用一对双荷子存在时电磁场的经典计算，推出 \eqref{eq:Line-Op-Locality}。

\subsubsection{例：$A_1$ 理论}\label{subsubsec:A1-labels}

本节讨论 $A_1$ 理论中线算子的标签。这些理论是由 $su(2)$ 规范群构造出的最一般的拉氏 $\CN=2$ 超共形场论。它们属于 \cite{Gaiotto:2009hg} 中记作 $\CS$ 的更大一类理论；背景与记号请参看该文。这些 $d=4$、$\CN=2$ 理论，是把一个部分扭转的 $ADE$ 型 $(0,2)$ 六维理论放在带穿刺点装饰的黎曼面 $C$ 上紧化而定义的。这类理论非凡的 $S$ 对偶性质在 \cite{Gaiotto:2009we} 中得到阐明。$\CS$ 类中的典型理论容许若干种 $S$ 对偶描述，即作为非阿贝尔规范理论弱耦合到非拉氏物质理论。我们预期，对非阿贝尔规范群存在 Wilson–'t Hooft 型的半 BPS 线算子；但非拉氏物质理论中可用的线算子该如何分类，我们目前尚不知道。$A_1$ 理论是一个有用的例子：它们与 $\CS$ 类理论共享许多有趣性质，却在所有 $S$ 对偶标架中都有拉氏量描述。

一个热身例子是具有 $\CN=4$ 超对称的 $su(2)$ 规范理论 \cite{Kapustin:2006pk}。此时只有一个余根 $H$，故 $\Lambda_{cr} = H \IZ$；也只有一个根 $\alpha$，故 $\Lambda_r = \alpha \IZ$。对偶配对为 $\langle \alpha, H\rangle = 2$，于是 $\Lambda_{wt} = \half \alpha \IZ$，$\Lambda_{mw} = \half H \IZ$。传统上分别选取同构 $\Lambda_{wt} \cong \IZ$ 与 $\Lambda_{mw} \cong \IZ$ 来定义电荷与磁荷；于是我们把 $(P,Q) = (\frac{p}{2} H, \frac{q}{2}\alpha)$ 等同于一对整数 $(p,q)\in \IZ\times \IZ$。相容性条件 \eqref{eq:Line-Op-Locality} 变成 $pq'-p'q= 0 \mod 2$。在这一同构下，若 $G = SU(2)$，则 $\Lambda_G \times \Lambda_G^* = 2\IZ \times \IZ$，即磁荷为偶数、电荷可为任意整数；而若 $G=SO(3)$，则 $\Lambda_G \times \Lambda_G^* =  \IZ \times 2\IZ$，即磁荷可为任意整数、而电荷为偶数。两者由 $S$ 对偶相联系。然而还有第三族相容的线算子集合，它不具有 $\CL_G$ 的形式，而是把 $SO(3)$ 理论的 $\theta$ 角平移半个周期得到的。\footnote{若用 $\frac{\theta}{8 \pi^2} \int_M {\Tr}_{2} F\wedge F$ 来归一化 $SO(3)$ 理论的 $\theta$ 角，则在一般四维流形 $M$ 上 $\theta$ 的周期是 $8 \pi$；但在自旋四维流形上是 $4\pi$。当然，$\IR^4$ 是自旋的。在这个意义上，我们平移了半个周期。}它由满足 $p=q \mod 2$ 的那些 $(p,q)$ 构成。

现在我们回到亏格 $g$、$n$ 个穿刺点的曲线 $C$ 上的 $A_1$ 理论。回忆：一个弱耦合极限由选取一个裤分解来指定，该分解与分离曲线 $c_i$（$i=1,\dots, 3g-3+n$）相伴。规范群的李代数为 $\lieg = su(2)^{\oplus 3g-3 + n}$；对每个由三条曲线 $(c_{i_1}, c_{i_2},c_{i_3})$ 围成的 trinion（三孔球面），都有一个超多重态 $\Phi_{i_1i_2i_3}$。这些超多重态处于表示 $2_{i_1}\otimes 2_{i_2} \otimes 2_{i_3}$ 中，并带一个实性条件（总共 $16$ 个独立实标量场）\cite{Gaiotto:2009we}。

首先注意：给定一个弱耦合尖点（cusp），可以有若干个可能的规范群。令 $\tilde G= SU(2)^{3g-3+n}$。此时容许的规范群形如 $\tilde G/\CZ$，其中 $\CZ$ 是中心 $Z(\tilde G)$ 的任一子群，只要它在所有 $2g-2+n$ 个超多重态上平凡地作用。为更具体地表达这一点，我们可以对每个 $i$ 配一个数 $\epsilon(c_i)\in \IZ_2$，从而固定一个同构 $Z(\tilde G) \cong \IZ_2^{3g-3+n}$。于是，在 trinion 上平凡作用的子群 $\CZ_{\rm max}$，由作用在向量 $\vec \epsilon$ 上的 $2g-2+n$ 个条件定义。这些关系之间还满足一个关系，因此 $\CZ_{\rm max}$ 由 $g$ 个元素生成。固定某个同构 $\CZ_{\rm max} \cong \IZ_2^g$，并记 $G_{\rm max}:= \tilde G/\CZ_{\rm max}$。在由分离曲线 $c_i$ 定义的尖点处，$A_1$ 理论的容许规范群与子群 $\CZ \subset \CZ_{\rm max}$ 一一对应。（顺带一提，这样的子群的个数是 Galois 数
\begin{equation}
\CG_g = 1+ \sum_{k=1}^g \frac{ (2^g-1)(2^{g-1}-1) \cdots (2^{g-k+1}-1)}{(2^k-1)(2^{k-1}-1)\cdots (2-1)}.
\end{equation}
当 $g$ 很大时 $ \CG_g \sim 2^{\frac{1}{4}g^2 }$，所以选择可以相当多。）

现在，对 $\lieg = su(2)^{\oplus 3g-3+n}$，我们把上面关于 $su(2)$ 的同构延拓为同构 $\Lambda_{mw} \times \Lambda_{wt} \cong \IZ^{3g-3+n} \times \IZ^{3g-3+n}$。相容的线算子集合，将由极大 $\CW$ 不变的向量族 $(\vec p, \vec q) \subset \Lambda_{G_{\rm max}} \times \Lambda_{wt}$ 构造，这些向量满足：对所有 $(\vec p, \vec q)$ 与 $(\vec p', \vec q')$ 都有 $\vec p\cdot \vec q' - \vec p' \cdot \vec q = 0 \mod 2$。向量 $\vec p$ 指定 $\IR^3-\{ 0 \} $ 上的一个 $SO(3)^{3g-3+n}$ 丛；限制 $\vec p \in \Lambda_{G_{\rm max}}$ 的出现，是因为超多重态场 $\Phi_{abc}$ 必须是伴随丛在表示 $2_a \otimes 2_b \otimes 2_c$ 中的截面。我们取 $p_i \geq 0$，并在 $p_i=0$ 时取 $q_i \geq 0$，以此固定外尔群作用的一个基本域。

Drukker、Morrison 与 Okuda 的一个漂亮观察 \cite{Drukker:2009tz} 是：$A_1$ 理论中单线算子的分类，与 $C$ 上非自交曲线之同痕类的 Dehn–Thurston 分类密切相关。回忆：对这种曲线 $c$，可以定义 $p_i = c \# c_i$ 与 $q_i = {\rm Twist}_{c_i}(c)$。这里 $\# $ 是同伦相交数，定义为相交数之和（取正），并对同痕类取极小。关于 twist 的细节见 \cite{Drukker:2009tz}。（下文中 $c$ 有时指一条曲线，有时指它的同痕类，视上下文而定。）相当漂亮的是，对每个 trinion 都有 $p_a + p_b + p_c = 0 \mod 2$，而这正是 $\vec p \in \Lambda_{G_{\rm max}}$ 的条件。因此我们可以定义一个一一到上的映射，记作 $\tau_{DMO}$，从 $C$ 上所有非自交路径的同痕类之集合，到 $(\Lambda_{G_{\rm max}} \times \Lambda_{wt})/\CW$。

那么，在这些理论中能定义哪些相容的线算子集合 $\CL$？在对应 $\tau_{DMO}$ 之下，条件 \eqref{eq:Line-Op-Locality} 等价于条件
\begin{equation}\label{eq:max-even}
c \# c' = 0 \mod 2.
\end{equation}
若（同痕类的）路径集合 $\CP$ 中任意一对都满足 \eqref{eq:max-even}，我们称 $\CP$ 为\emph{偶}的；若 $\CP$ 不真包含于任何更大的偶集合之中，则称其为\emph{极大偶}的。于是，相容的线算子集合 $\CL$ 定义了 $C$ 上一个极大偶路径集合。特别地，若对某个子群 $\CZ \subset \CZ_{\rm max}$ 取 $G = \tilde G/\CZ$，则 $(\Lambda_G \times\Lambda_G^*)/\CW$ 决定一个极大偶路径集合，记作 $\CP_G$。反之，给定 $\CP$，我们可以把它投影到磁荷，并重构出与 \eqref{eq:label-seq} 一样的正合列。

$S$ 对偶群，即模群，会把规范群变掉，正如上面提到的 $\CN=4$ 例子那样。用模二同调可以给出 $\CZ$ 如何变化的显式公式。这补足了 \cite{Drukker:2009tz} 中关于 $S$ 对偶对线算子作用的讨论。
